\section{Integral turning data from a rotation system}
\label{sec:turning}

\subsection{Darts, faces, and smoothing conventions}

Suppose first that $n>0$. A validated PD diagram is a connected four-valent
plane multigraph with $n$ vertices and $2n$ edges. Loops and repeated edges
are allowed. At crossing $v$, its four darts are numbered $(v,0)$ through
$(v,3)$ in counterclockwise order. Write $\sigma(v,j)=(v,j+1\bmod4)$ and
let $\alpha$ exchange the two darts of each edge. The cycles of
\begin{equation}\label{eq:face-permutation}
 \phi=\sigma\alpha
\end{equation}
are the oriented face walks. There are $n+2$ of them, and their total length
is $4n$. A face walk may repeat vertices and edges.

The two smoothings pair slots as follows:
\[
 S_0=\{\{0,1\},\{2,3\}\},\qquad
 S_1=\{\{0,3\},\{1,2\}\}.
\]
Their bracket weights are $A$ and $A^{-1}$, respectively. We use the
unnormalized bracket $\mathcal B_D(A)$, in which every closed circle has
weight $\delta=-A^2-A^{-2}$. Thus
\begin{equation}\label{eq:bracket-state-sum}
 \mathcal B_D(A)=\sum_{s\in\{0,1\}^n}
      A^{n-2|s|}\delta^{c(s)},
 \qquad
 V_D(A^{-4})=(-A^3)^{-\operatorname{wr}(D)}
                   \frac{\mathcal B_D(A)}{\delta}.
\end{equation}
Here $c(s)$ counts the smoothing circles. These conventions agree with the
independent Laurent-polynomial oracle supplied by the baseline.

At a smoothing arc traversed from incoming slot $j$ to outgoing slot $k$,
define its quarter-turn by
\begin{equation}\label{eq:corner-turn}
 \tau(j,k)=
 \begin{cases}
 -1,&k=j+1\pmod4,\\
 +1,&k=j-1\pmod4.
 \end{cases}
\end{equation}
The incoming tangent points into the crossing, which explains the minus
sign in the first case. In particular, each corner of a face walk in
\eqref{eq:face-permutation} has quarter-turn $-1$.

\begin{figure}[ht]
\centering
\begin{tikzpicture}[scale=0.85,>=Latex]
 \draw[fill=pale,draw=navy] (0,0) circle (1.0);
 \foreach \a/\s in {0/0,90/1,180/2,270/3}{
  \draw[navy] (\a:0.98)--(\a:1.65);
  \node at (\a:1.92) {$\s$};
 }
 \draw[very thick,teal,->] (0,1.4)--(0,0.72)
  .. controls (0,0.3) and (0.3,0) .. (0.72,0)--(1.4,0);
 \node[align=center] at (0,-2.65) {incoming $1$ to outgoing $0$\\$\tau(1,0)=+1$};
 \begin{scope}[xshift=6.0cm]
 \draw[fill=pale,draw=navy] (0,0) circle (1.0);
 \foreach \a/\s in {0/0,90/1,180/2,270/3}{
  \draw[navy] (\a:0.98)--(\a:1.65);
  \node at (\a:1.92) {$\s$};
 }
 \draw[very thick,teal,->] (1.4,0)--(0.72,0)
  .. controls (0.3,0) and (0,0.3) .. (0,0.72)--(0,1.4);
 \node[align=center] at (0,-2.65) {incoming $0$ to outgoing $1$\\$\tau(0,1)=-1$};
 \end{scope}
\end{tikzpicture}
\caption{The local convention. Slot numbers increase counterclockwise,
but the entering tangent reverses the slot's outward direction.}
\label{fig:turn-convention}
\end{figure}

\subsection{A small system of face equations}

Choose one face $f_\infty$ as the outer face. A directed edge-turn cochain
is a family of integers $b_d$, one for every dart, satisfying
$b_{\alpha(d)}=-b_d$. We require
\begin{equation}\label{eq:turn-equations}
 \sum_{d\in f} b_d=
 \begin{cases}
 |f|-4,&f\ne f_\infty,\\
 |f|+4,&f=f_\infty.
 \end{cases}
\end{equation}
These equations contain the total edge turning along each face; the
$|f|$ local corners contribute $-|f|$ separately.

\begin{theorem}[Constructive integral turning cochain]
\label{thm:turn-construction}
For every connected spherical four-valent rotation system with $n>0$, the
equations \eqref{eq:turn-equations} have an integer solution supported on
the edges dual to a dual spanning tree. Such a solution can be constructed
and verified in polynomial bit time and encoded in $O(n\log(n+2))$ bits.
For the construction used here,
\begin{equation}\label{eq:cochain-bounds}
 \max_d |b_d|\le4n+4,
 \qquad
 L_b:=\sum_{e}|b_e|\le4(n+1)^2.
\end{equation}
The sum over edges chooses either orientation once.
\end{theorem}

\begin{proof}
Put $h_f=|f|-4+8\mathbf1_{f=f_\infty}$. Euler's formula gives
\[
 \sum_f h_f=4n-4(n+2)+8=0.
\]
Choose a spanning tree of the connected dual graph, rooted at $f_\infty$.
For a nonroot face $f$, let $d_f$ be the dart on its parent edge belonging
to $f$. Assign $b_{d_f}$ the total demand of the subtree rooted at $f$ and
assign its negative to the other dart. Assign zero on all remaining edges.
At $f$, its parent-edge contribution is its subtree demand; the
child-edge contributions subtract the children's subtree demands.
Their sum is $h_f$. The root equation follows from total demand zero.
This also handles a primal bridge: its dual loop belongs to no spanning
tree and contributes zero twice with opposite signs.

The total positive demand equals half the absolute demand sum, and
\[
 \sum_f |h_f|\le\sum_f(|f|+4)=8n+8.
\]
Every subtree sum therefore has absolute value at most $4n+4$.
There are $n+1$ dual-tree edges, proving \eqref{eq:cochain-bounds}.
The construction uses face walks, a tree traversal, and integer subtree
sums. Its verifier reconstructs the face walks, checks antisymmetry, and
checks every equation in \eqref{eq:turn-equations}; it does not repeat the
tree construction. All numbers used by the constructed certificate have
$O(\log(n+2))$ bits.
\end{proof}

\subsection{Why every smoothing circle has the right turning}

\begin{lemma}[Zero face integrals determine all closed integrals]
\label{lem:closed-cochain}
Let $c$ be an antisymmetric edge cochain on a connected graph cellularly
embedded in the sphere. If its integral around every face walk is zero,
then its integral around every closed graph walk is zero. The statement
holds over $\Z$ and over $\Z/4\Z$.
\end{lemma}

\begin{proof}
Choose a primal spanning tree. For each nontree edge, its fundamental
cycle bounds a union of faces, with signs determined by orientation.
The integral around that cycle is consequently zero. Define a vertex
potential by integrating $c$ along the tree from a root. The vanishing
fundamental-cycle integrals show that the difference of potentials across
every edge is $c$ on that edge. Potential differences telescope on every
closed walk. This argument uses no division and works over either ring.
Repeated edges on a face boundary cancel with their orientations.
\end{proof}

\begin{theorem}[Turning of all smoothing circles]
\label{thm:turning}
Let $b$ satisfy \eqref{eq:turn-equations}. For any smoothing state, orient
one of its circles $C$ arbitrarily. Sum $b_d$ along its directed traversals
of diagram edges and sum \eqref{eq:corner-turn} along its smoothing arcs.
The resulting integer is $+4$ or $-4$. Reversing $C$ negates the integer.
\end{theorem}

\begin{proof}
Take any smooth planar realization of the rotation system with the chosen
outer face. In small disjoint neighborhoods of its crossings, arrange
the four outward tangents in the cardinal directions indexed by the
slots. This is possible because their cyclic orders agree. Smooth the
remaining edges while keeping those endpoint tangents fixed.
Let $b^{\mathrm{geom}}_d$ be the total tangent turning along such an edge,
measured in quarter-turns. Endpoint tangent directions make this an integer.
It is antisymmetric under reversal.

A slightly regularized face boundary turns once clockwise for an inner
face and once counterclockwise for the outer face. Its local crossing
corners each contribute $-1$. Hence $b^{\mathrm{geom}}$ satisfies exactly
\eqref{eq:turn-equations}, including faces whose combinatorial walks
repeat vertices. By Lemma~\ref{lem:closed-cochain}, the difference
$b-b^{\mathrm{geom}}$ has zero integral on every closed graph walk.

The oriented smoothing circle is a simple smooth plane curve. Its total
tangent turning is $+4$ or $-4$ quarter-turns. Replacing the edge part of
that sum by $b$ leaves the total unchanged: after collapsing the small
crossing neighborhoods, the circle determines a closed graph walk.
It may visit one crossing twice, but the lemma explicitly permits this.
The local smoothing turns are precisely \eqref{eq:corner-turn}.
Reversal negates both the edge and corner contributions.
\end{proof}

This proof uses a geometric drawing only as an existence argument. The
algorithm constructs and verifies integers from the rotation system.
No coordinate drawing, floating-point angle, or numerical winding test
appears in the computation.

\begin{corollary}[Gauge independence]
\label{cor:gauge}
For a fixed outer face, any two solutions of the face equations give the
same total turn on every smoothing circle. Different outer-face choices
also yield the same unoriented loop evaluation $\rho^4+\rho^{-4}$.
\end{corollary}

\begin{proof}
The first assertion is Lemma~\ref{lem:closed-cochain}. For the second,
Theorem~\ref{thm:turning} still gives the two opposite values $\pm4$;
summing over both orientations removes their order.
\end{proof}

\begin{remark}[The sphere hypothesis]
On a higher-genus surface, vanishing face integrals need not imply
vanishing integrals on essential cycles. The construction here therefore
requires validated classical diagrams. It is not a Jones algorithm for
arbitrary virtual diagrams with no additional homological data.
\end{remark}
